\documentclass{article}
\usepackage[utf8]{inputenc}
\usepackage[ukrainian]{babel}
\usepackage[T2A]{fontenc}
\usepackage{hyperref}
\usepackage{enumitem}

\title{Історія розвитку N2O: високосумісна лінія протоколів для розподілених систем}
\author{Максим Сохацький \\ ДП "Інформаційні судові системи"}
\date{\today}

\begin{document}

\maketitle

\begin{abstract}
Стаття досліджує історію бібліотеки N2O від SYNRC --- вбудованої бібліотеки циклу протоколів
повідомлень для WebSocket, MQTT, TCP та інших серверів. Особлива увага приділяється високому
рівню сумісності лінії N2O та єдиній ключовій біфуркаційній точці в її еволюції: переходу від
чистого WebSocket-тракту до абстрактних трактів, що дозволило міграцію на інші протоколи,
зокрема MQTT.
\end{abstract}

\tableofcontents

\section{Introduction}
N2O (Nitrogen 2.0 Optimized) --- це легковагова вбудована бібліотека для мов Erlang і Elixir для обробки
повідомлень у реальному часі на основі веб-фреймворку Nitrogen від Расті Клопхауза.
Вона забезпечує високопродуктивний релей протоколів, керування процесами, віртуальні кільця вузлів,
сесії та уніфікований API для зовнішніх сервісів черг і кешування.

Проєкт містить 700 рядків коду та сумісний з різними хостами (Cowboy, Bandit, EMQ) і шинами повідомлень (gproc, syn, pg2).
N2O демонструє принцип мінімалізму та еволюційного розвитку, де більшість змін зберігають зворотну сумісність,
а архітектура дозволяє просту міграцію між транспортними протоколами.

\section{SYNRC N2O History}

\subsection{2012--2014: Початок та фокус на WebSocket}
N2O виник у рамках екосистеми Synrc як потужний Erlang web-фреймворк, натхненний Nitrogen.
Початкова версія орієнтувалася на чистий WebSocket-тракт з Cowboy. Основний акцент робився
на простоті, продуктивності та інтеграції з Erlang OTP. У цей період сформувалася базова
архітектура: n2o\_proto, процеси (n2o\_pi) та підтримка форматів (BERT, JSON). Лінія
демонструвала високу сумісність з існуючими Erlang-інструментами.

\subsection{2015--2016: Розширення функціональності}
Розвиток Nitro (шаблонізатор) та інтеграція з KVS (абстрактним шаром БД). WebSocket
залишався основним трактом. Проєкт набув популярності в спільноті Erlang завдяки мінімалізму
та продуктивності реального часу. Підтримка різних pub/sub бекендів підкреслювала гнучкість.

\subsection{2017--2018: Абстракція трактів}
Ключовий момент еволюції: перехід від чистого WebSocket-тракту до абстрактних протоколів (n2o\_proto).
Це єдина значна біфуркація в історії лінії N2O. Абстракція дозволила безболісну міграцію на інші протоколи,
зокрема MQTT (через emqttd). Уніфікація версій WebSocket і MQTT (5.11.1) стала кульмінацією цього переходу.
Після цієї точки N2O зберіг високу зворотну сумісність, демонструючи архітектурну стійкість.

\subsection{2019--2021: Уніфікація та зрілість}
Підтримка Elixir (Mix), оновлення до OTP 21+, покращення обробки помилок, мінімалізація коду.
Додавання підтримки arraybuffer у JS-клієнті. Релізи 2020 року (до 6.x) закріпили статус N2O
як універсальної бібліотеки для WebSocket/MQTT/TCP. Проєкт використовувався в
enterprise-додатках (банківські системи, IoT, ERP).

\subsection{2022--2026: Підтримка та стабільність}
Продовження розвитку з фокусом на стабільність, нові хости (Bandit) та інтеграції.
Репозиторій накопичив понад 2300 комітів. N2O залишається прикладом довготривалої,
сумісної лінії з мінімальними breaking changes.

\subsection{2026--2030: Перехід на ASN.1 DER кодування}
Керуючись підвищеними вимогами до телекомунікаційних продуктів критичної інфраструктури
було прийнято рішення про поступову міграцію сервісів на ASN.1 визначені протоколи і
промисловий формат серіалізації який використовується в безпековій інфраструктурі X.509.

\section{Conclusion}
Історія N2O ілюструє успішну еволюцію мінімалістичного підходу в розробці розподілених систем.
Завдяки високому рівню сумісності та лише одній ключовій біфуркаційній точці (абстракції трактів),
лінія N2O дозволила органічний перехід від WebSocket-центричної архітектури до універсальної підтримки протоколів,
включаючи MQTT. Це робить N2O потужним прикладом для сучасних real-time систем, де простота,
продуктивність та майбутня міграція є критичними.

\begin{thebibliography}{9}
\bibitem{github} GitHub Repository: synrc/n2o. \url{https://github.com/synrc/n2o}
\bibitem{readme} README N2O. Офіційна документація.
\bibitem{history} HISTORY.md, synrc/n2o.
\bibitem{medium} Namdak Tonpa. MQTT + N2O = Synrc Messaging Core. Medium, 2017.
\bibitem{pdf} Maxim Sokhatsky. N2O Presentation. Erlang Factory.
\end{thebibliography}

\end{document}
